\section{Uniform localization of the entire variational sublevel}
\label{sec:sc-sublevel}
\subsection{Actual records and the inherited comparison functional}
All lengths are in units of the triangular lattice spacing. Let
\[
 \Lambda=\{(a+b/2,\sqrt3b/2):a,b\in\mathbb Z\},\qquad
 R\in[9/20,47/100],\quad R_+=47/100.
\]
Fix any actual regular record of $m\ge1$ free flights and $m+1$ collision points $Q_0,\ldots,Q_m$. Its successive lattice centers are $O_0,\ldots,O_m$, with $O_i\ne O_{i+1}$. Write
\[
 Q_i=O_i+Rn_i,\quad |n_i|=1,\quad
 Q_{i+1}-Q_i=\tau_i v_i,\quad |v_i|=1.
\]
Regular means genuine exterior, first-hit, nontangent specular collisions at all these points and no extra collision on any open flight. Nonconsecutive repeated centers are permitted. At interior collisions let
\begin{equation}\label{eq:sc-reflection}
 c_i=-v_{i-1}\cdot n_i=v_i\cdot n_i>0,
 \qquad v_i-v_{i-1}=2c_i n_i\quad(1\le i<m).
\end{equation}
The endpoint cosines are $c_0=v_0\cdot n_0>0$ and $c_m=-v_{m-1}\cdot n_m>0$. Define
\begin{equation}\label{eq:sc-defect}
 L=\sum_{i=0}^{m-1}\tau_i,\qquad E=R(2-c_0-c_m),
 \qquad 0\le E\le2R\le47/50<1.
\end{equation}
The original endpoint momenta obey $p_0^2+c_0^2=p_m^2+c_m^2=1$, so $E=R[p_0^2/(1+c_0)+p_m^2/(1+c_m)]$.

For this fixed ordered word, the comparison domain and length are
\[
 \mathcal K=\prod_{i=0}^m\overline B(O_i,R),\qquad
 A(Y)=\sum_{i=0}^{m-1}|Y_{i+1}-Y_i|.
\]
Repeated centers give separate factors. This is the native disk-product comparison, not a substitute physical dynamics. Define its full sublevel
\[
 \mathcal S(Q)=\{Y\in\mathcal K:A(Y)\le L\}.
\]
It is nonempty and compact. Every global minimum of $A$ lies in it, whether or not the minimum is physical. For $Y\in\mathcal K$ put
\[
 \delta_i=Y_i-Q_i,\qquad
 Y_{i+1}-Y_i=\ell_i\widehat v_i,\qquad e_i=\widehat v_i-v_i.
\]
Consecutive disks are separated by $\tau_-=1-2R_+=3/50$, so every $\ell_i\ge\tau_->0$. The retained geometric facts, reconstructed in Appendix~\ref{sec:sc-geometry}, are
\begin{equation}\label{eq:sc-geometry-input}
 \tau_i\le H:=28/5,\qquad
 \sum_{i=a+1}^{a+93}c_i\ge\gamma:=1/4480
 \quad\text{whenever }0\le a,\ a+94\le m.
\end{equation}
Thus every complete 94-flight block has an interior index with
$c_i\ge c_{\rm a}:=1/416640$. No such condition is imposed on all individual incidences.

\subsection{The complete positive gap, not a coordinatewise singular estimate}
The next calculation extends the inherited endpoint-gap bookkeeping from minima to the entire sublevel. The general disk duality and the endpoint defect themselves were already established in module 297.
\begin{sclemma}[Full sublevel energy and contact budget]\label{lem:sc-gap}
For every $Y\in\mathcal S(Q)$,
\begin{align}
 \sum_{i=0}^{m-1}\ell_i|e_i|^2&\le2E,\label{eq:sc-flight-energy}\\
 \sum_{i=1}^{m-1}c_i|\delta_i|^2&\le RE,\label{eq:sc-contact-energy}\\
 |\delta_0|,|\delta_m|&\le2\sqrt{2RE}.\label{eq:sc-endpoints}
\end{align}
These statements hold for interior or inactive comparison vertices as well as boundary vertices.
\end{sclemma}
\begin{proof}
Set $d_0=-v_0$, $d_i=v_{i-1}-v_i=-2c_i n_i$ for $1\le i<m$, and $d_m=v_{m-1}$. Telescoping, including both endpoints exactly once, gives
\[
 \sum_{i=0}^{m-1}v_i\cdot(Y_{i+1}-Y_i)=\sum_{i=0}^m d_i\cdot Y_i.
\]
Consequently, with $D=\sum_i(d_i\cdot O_i-R|d_i|)$,
\begin{equation}\label{eq:sc-positive-gap}
 A(Y)-D=\frac12\sum_{i=0}^{m-1}\ell_i|e_i|^2+
 \sum_{i=0}^m\bigl[d_i\cdot(Y_i-O_i)+R|d_i|\bigr].
\end{equation}
Every term is nonnegative: the first uses two unit vectors, and the second is the difference from the minimum of a linear functional on its actual disk. At $Q$ every flight and interior contact term is zero; the endpoint terms are $R(1-c_0)$ and $R(1-c_m)$. Thus $L-D=E$. Since $A(Y)\le L$, the whole right-hand side of \eqref{eq:sc-positive-gap} is at most $E$, proving \eqref{eq:sc-flight-energy}.

For any $z\in\overline B(0,R)$ and unit $n$, the exact identity
\begin{equation}\label{eq:sc-disk-curvature}
 2R(R-n\cdot z)=|z-Rn|^2+R^2-|z|^2
\end{equation}
implies $|z-Rn|^2\le2R(R-n\cdot z)$. An interior contact term in \eqref{eq:sc-positive-gap} equals $2c_i(R-n_i\cdot(Y_i-O_i))$ and is therefore at least $c_i|\delta_i|^2/R$. Summing proves \eqref{eq:sc-contact-energy}, without dividing by any individual $c_i$.

At the first endpoint let $Z_0=O_0+Rv_0$. The endpoint term gives $|Y_0-Z_0|^2\le2RE$ by \eqref{eq:sc-disk-curvature}, while $|Q_0-Z_0|^2=2R^2(1-c_0)\le2RE$. The triangle inequality proves the first bound in \eqref{eq:sc-endpoints}. At the final endpoint use $Z_m=O_m-Rv_{m-1}$; the signs give the same bound.
\end{proof}

\begin{sclemma}[Transverse telescoping with the turning work retained]\label{lem:sc-telescoping}
Put $w_i=v_i^\perp$ and $r_i=\ell_i e_i$. For $0\le j<i\le m$,
\begin{equation}\label{eq:sc-transverse}
 w_{i-1}\cdot\delta_i=w_j\cdot\delta_j+
 \sum_{l=j}^{i-1}w_l\cdot r_l+
 \sum_{l=j+1}^{i-1}(w_l-w_{l-1})\cdot\delta_l.
\end{equation}
If $i-j\le94$ and $Y\in\mathcal S(Q)$, the absolute values of the last two sums are respectively bounded by $38\sqrt E$ and $14\sqrt E$.
\end{sclemma}
\begin{proof}
The exact displacement equation is
\[
 \delta_{l+1}-\delta_l=(\ell_l-\tau_l)v_l+r_l.
\]
Its projection on $w_l$ removes only its longitudinal term. Summing these projected equations and collecting coefficients of each $\delta_l$ yields \eqref{eq:sc-transverse}. The last sum is present because the transverse frame turns. It cannot in general be omitted.

All comparison vertices obey $|\delta_l|\le2R$, so
\[
 \ell_l\le\tau_l+|\delta_l|+|\delta_{l+1}|
 \le H+4R_+=187/25=:T.
\]
Weighted Cauchy--Schwarz and \eqref{eq:sc-flight-energy} give
\[
 \sum_{l=j}^{i-1}|r_l|\le
 \left(\sum_{l=j}^{i-1}\ell_l\right)^{1/2}
 \left(\sum_{l=j}^{i-1}\ell_l|e_l|^2\right)^{1/2}
 \le\sqrt{2\cdot94\cdot T E}<38\sqrt E.
\]
The rotation is an isometry, so \eqref{eq:sc-reflection} gives $|w_l-w_{l-1}|=2c_l$. Therefore \eqref{eq:sc-contact-energy} gives
\[
 \sum_{l=j+1}^{i-1}|(w_l-w_{l-1})\cdot\delta_l|
 \le2\left(\sum_{l=j+1}^{i-1}c_l\right)^{1/2}
          \left(\sum_{l=j+1}^{i-1}c_l|\delta_l|^2\right)^{1/2}
 \le2\sqrt{94R_+E}<14\sqrt E.
\]
Only actual reflection cosines enter the turning sum. In particular no dynamics is ascribed to the comparison polygon.
\end{proof}

\begin{sctheorem}[Collision-uniform sublevel localization without an incidence cutoff]\label{thm:sc-quarter}
For every actual regular record in the original radius family, every $m\ge1$, and every $Y\in\mathcal S(Q)$,
\begin{equation}\label{eq:sc-quarter}
 \max_{0\le i\le m}|Y_i-Q_i|\le22\,E^{1/4}.
\end{equation}
The constant is independent of the word, collision count, section/return/displacement labels and every interior incidence. At $E=0$ the sublevel is the singleton $\{Q\}$. No optimality of the exponent or nontriviality for large $E$ is asserted.
\end{sctheorem}
\begin{proof}
If $E=0$, \eqref{eq:sc-contact-energy} and positivity of the actual interior cosines force all interior displacements to vanish, and \eqref{eq:sc-endpoints} handles the endpoints. Assume $0<E<1$.

Call an interior index an anchor if $c_i\ge c_{\rm a}$. At such an index,
\[
 |\delta_i|\le\sqrt{R_+/c_{\rm a}}\sqrt E<443\sqrt E.
\]
Together with $|\delta_i|\le2R_+$ this gives
$|\delta_i|^2<2R_+443\sqrt E<22^2\sqrt E$.
At either endpoint \eqref{eq:sc-endpoints} gives
$|\delta_i|^2\le8R_+E<4\sqrt E<22^2\sqrt E$.

Consider a nonanchor interior index $i$. Let $j<i$ be the last anchor before it, or $j=0$ if none exists. Necessarily $i-j<94$: otherwise the 94-flight block starting at $j$ has by \eqref{eq:sc-geometry-input} an anchor strictly between $j$ and $i$, a contradiction. The initial displacement in \eqref{eq:sc-transverse} is at most $443\sqrt E$, also when $j=0$ by \eqref{eq:sc-endpoints}. Lemma~\ref{lem:sc-telescoping} therefore yields
\[
 |w_{i-1}\cdot\delta_i|\le(443+38+14)\sqrt E=495\sqrt E.
\]
Because $n_i\cdot v_{i-1}=-c_i$ and $|n_i\cdot w_{i-1}|\le1$,
\[
 |n_i\cdot\delta_i|\le c_i|\delta_i|+|w_{i-1}\cdot\delta_i|.
\]
Crucially, the first term needs no inverse incidence: by \eqref{eq:sc-contact-energy},
$c_i|\delta_i|\le\sqrt{c_i}\sqrt{c_i|\delta_i|^2}\le\sqrt{R_+E}<\sqrt E$.
Thus $|n_i\cdot\delta_i|\le496\sqrt E$. Since $Y_i$ lies in the disk and $Q_i$ on its boundary, \eqref{eq:sc-disk-curvature} becomes
\[
 |\delta_i|^2\le-2R n_i\cdot\delta_i
 \le2R_+496\sqrt E=\frac{11656}{25}\sqrt E
 <484\sqrt E.
\]
Taking square roots proves \eqref{eq:sc-quarter}. The argument treats short records and the first incomplete block without inserting an artificial collision. It uses at most one bounded physical block for each index, not a sum over the full record or over its words.
\end{proof}

\begin{scremark}[What is new relative to the complete predecessor]
The earlier pointwise contact certificate had a factor $c_i^{-1/2}$. Here the contact budget is used \emph{before} division: its weighted sum pays the turning work in \eqref{eq:sc-transverse}, and disk curvature converts the resulting normal control into an unweighted position bound. The dense-anchor input is inherited real geometry, not a fresh probability hypothesis. The result concerns every point of a continuous compact sublevel, not a list of sampled minima. It neither constructs an orbit through a singularity nor controls forward-time derivatives at grazing.
\end{scremark}
